%% kapa: exam=1 round=25 year=2014 period=1 form=A subject=civil_law subject_ko=민법 seq=28 number=28
%% source: Q-Net 감정평가사 기출문제 (cst00309) artlSeq=5101207 file=extracted/1008/5101207/2014년 제25회 감정평가사 1차시험 기출문제/1교시(A형).hwp
%% conversion: kordoc 4.24.51 -> tools/gen1.py (auto); review=proofread-sonnet; flags=none
\begin{question}{28}
부동산의 점유취득시효에 관한 설명으로 옳은 것은? (다툼이 있으면 판례에 의함)
\begin{choices}
  \item 자기소유의 부동산은 취득시효의 대상이 될 수 없다.
  \item 시효기간이 진행되던 중에 등기부상의 소유자가 변경된 경우, 시효완성자는 시효완성 당시의 소유자에 대하여 시효완성의 효과를 주장할 수 있다.
  \item 취득시효를 주장하는 자는 시효기간의 만료로 즉시 소유권을 취득한다.
  \item 취득시효를 주장하는 자는 반드시 실제로 점유를 개시한 때를 시효의 기산점으로 삼아야 한다.
  \item 취득시효가 완성되었더라도 등기를 하지 않고 있는 사이에 제3자가 그 부동산에 관한 소유권 이전등기를 마친 경우, 시효완성자는 원칙적으로 그 제3자에 대하여 시효완성의 효과를 주장할 수 있다.
\end{choices}
\end{question}
\answer{2}
